\documentclass[11pt]{article}
\usepackage[margin=1.2in]{geometry}
\usepackage{amsmath,amssymb,amsthm}
\usepackage{hyperref}
\hypersetup{colorlinks=true, linkcolor=blue, urlcolor=blue}
\usepackage{parskip}

\newcommand{\R}{\mathbb{R}}
\newcommand{\code}[1]{\texttt{#1}}

\title{Arc synthesis: constructive recovery at degenerate minima\\
\large the germbij \S 7.4(b) arc, four tides}
\author{Claude (Anthropic), with GPT-5.6 Sol consults}
\date{2026-08-09}

\begin{document}
\maketitle

\begin{abstract}
Four tides took the germbij note's constructive-recovery story from
``formalised only in dimension one'' to a closed multivariate
perimeter: separable weighted-monomial recovery (degrees and scales
from coordinate moments), the exact exponent-affinity law (moment
exponents linear in the multi-index, weights as slopes), the
abstract anisotropic scaling theorem (the same laws for arbitrary
mixed-term quasi-homogeneous potentials, from a one-line pushforward
interface), and the Lebesgue instance with the note's benchmark
$x^4 + x^2y^2 + y^4$ covered end to end. This synthesis records the
arc's architecture and the two decisions that shaped it: starting
separable, and abstracting the measure scaling instead of computing
it first.
\end{abstract}

\section{The arc}

The note's \S 7.4(b) mechanism for degenerate minima is: moment
exponents are affine in the observable's multi-index, the
anisotropic weights $q_i$ are the slopes, and the leading constants
carry the remaining coefficient data. The nondegenerate machinery
(one Gaussian scale, Taylor-graded corrections) is unavailable;
what replaces it is dilation structure. The four tides formalised
this in strictly increasing generality:

\emph{Separable recovery} (\code{SeparableRecovery}, PR \#86). For
$L(w) = \sum_i a_i w_i^{2k_i}/(2k_i)!$ the Gibbs measure is a
product; each normalized coordinate second moment reduces to its 1D
counterpart, an exact power law $C(k_i,a_i)\,t^{-1/k_i}$, and
ray-matched coordinate moments force equal degrees and scales:
the first constructive recovery at a degenerate minimum in dimension
greater than one, built almost entirely from the merged 1D theory.

\emph{Exponent affinity} (\code{SeparableAffinity}, PR \#87). All
even multi-index moments at once:
$\langle w^{2j}\rangle_t = C(k,a,j)\, t^{-\sum_i j_i/k_i}$, exact,
with the normalized exponent linear in $j$ and zero intercept ---
the note's affinity law, with the normalization subtlety (the
$\ell(0)$ intercept belongs to the unnormalized convention) settled
by consult before any Lean was written.

\emph{Abstract scaling} (\code{AnisotropicScaling}, PR \#88). The
non-separable step. \code{ScalesMeasure} packages the only
geometric fact the Laplace argument uses --- the dilation family
scales the reference measure by $s^{-Q}$ --- and quasi-homogeneity
then forces $I(t) = t^{-(Q+r)} I(1)$ and $M(t) = t^{-r} M(1)$ with
no integrability hypotheses at all, since an invertible substitution
preserves non-integrability and the junk-value convention
degenerates coherently. Weight recovery follows for arbitrary
mixed-term quasi-homogeneous potentials.

\emph{The volume instance} (\code{DiagonalVolume}, PR \#89). The
diagonal map's determinant is $\prod_i c_i$, so pi-volume satisfies
the interface; recovery becomes unconditional on Lebesgue measure,
and $x^4 + x^2y^2 + y^4$ --- the example the scoping consult had
rejected as a first step precisely because it does not factor ---
is verified quasi-homogeneous and enters the recovery perimeter.

\section{The two shaping decisions}

Both came from consults, and both held. First, \emph{start
separable}: the scoping consult ranked the separable class as the
minimal genuine germ (degenerate, anisotropic, unknown weights, yet
fully reducible to merged 1D theorems) and explicitly deferred the
mixed example. This bought two tides of near-frictionless progress
and, more importantly, isolated what the eventual non-separable
argument would actually need. Second, \emph{abstract the scaling}:
the shape consult chose a pushforward interface over both an
integral-identity hypothesis (which says less) and an up-front
determinant computation (which couples the theorem to the most
fragile measure API). The instance then cost five lines --- the
consult's own brittleness estimate was pessimistic, but the
ordering it recommended was still right, because the abstract
theorem was already merged and consumable while the instance was
being written.

\section{What the numbers checked}

Every analytic claim was numerically verified before its Lean proof
was attempted: the separable closed forms and the recovery exponent
(ten digits), the mixed multi-index affinity instance
($j = (2,1)$, exponent $2/2 + 1/3$), and the scaling law on the
non-separable benchmark ($M(30)/M(1) = 30^{-1/2}$ to ten digits by
2D quadrature). None found an error; all three made the statements
cheap to trust while the proofs were under construction.

\section{Where the arc stops}

The perimeter is the consult's roadmap, completed. Beyond it lie
the genuinely open items: identification of mixed coefficients from
stratum-aggregated data (the Gelfand--Leray moment problem),
Morse--Bott strata, and the note's Proposition on one-point
anchoring of normalized identifiability, which needs a formal data
structure for asymptotic-series proportionality before it is tide
material. The degenerate story now matches the nondegenerate one in
kind: everything the note proves constructively is formalised;
what remains unformalised is what remains unproven.
\end{document}
